\documentclass[12pt,a4paper]{article}
\usepackage{iftex}
\ifPDFTeX
  \usepackage[T2A]{fontenc}
  \usepackage[utf8]{inputenc}
  \usepackage[english,russian]{babel}
\else
  \usepackage{fontspec}
  \setmainfont{DejaVu Serif}
  \setmonofont{DejaVu Sans Mono}
  \renewcommand{\contentsname}{Содержание}
  \renewcommand{\figurename}{Рис.}
  \renewcommand{\tablename}{Таблица}
\fi

\usepackage[a4paper,margin=25mm]{geometry}
\usepackage{amsmath,amssymb,amsthm}
\usepackage{graphicx}
\usepackage{booktabs}
\usepackage{enumitem}
\usepackage[hidelinks]{hyperref}

\newcommand{\Exp}{\mathsf{M}}
\newcommand{\Var}{\mathsf{D}}
\newcommand{\PP}{\mathsf{P}}
\newcommand{\dens}{\mathrm{p}}
\newcommand{\Normal}{\mathcal{N}}
\newcommand{\R}{\mathbb{R}}
\newcommand{\xs}[1]{x^{(#1)}}
\newcommand{\ys}[1]{y^{(#1)}}
\newcommand{\norm}[1]{\left\lVert #1 \right\rVert}
\newcommand{\abs}[1]{\left\lvert #1 \right\rvert}
\newcommand{\dd}{\,\mathrm{d}}
\newcommand{\est}[1]{\tilde{#1}}
\newcommand{\zad}[1]{\medskip\noindent\textbf{Задание #1.}\quad}
\newcommand{\prov}{\smallskip\noindent\textit{Проверка:}\ }

\theoremstyle{definition}
\newtheorem{definition}{Определение}
\newtheorem{statement}{Утверждение}
\theoremstyle{plain}
\newtheorem{lemma}{Лемма}

% Список с русскими буквенными метками
\newenvironment{ruslist}
  {\begin{itemize}[label={},leftmargin=2.2em,itemsep=2pt,topsep=3pt]}
  {\end{itemize}}


\begin{document}

\begin{center}
  {\large\bfseries Задача 4. Нейронные сети}\\[1mm]
  Нейронные сети и машинное обучение\\[1mm]
  Задания 1--8
\end{center}

\section{Что нужно знать}

\subsection{Обозначения}

\begin{center}
\begin{tabular}{ll}
\toprule
$x \in \R^m$ & вектор признаков \\
$w$ & вектор весов (параметров) \\
$\langle w, x\rangle$ & скалярное произведение \\
$\theta$ & ступенчатая функция активации \\
$\eta$ & коэффициент обучения \\
$\mathcal{C}_1, \mathcal{C}_2$ & классы \\
$\mathcal{H}_1, \mathcal{H}_2$ & обучающие выборки этих классов \\
$\est w$ & разделяющий вектор весов \\
$w(n)$ & веса на $n$-й итерации \\
$W_l$, $b_l$ & матрица весов и вектор сдвигов $l$-го слоя \\
$\sigma$, $\tanh$ & логистический сигмоид и гиперболический тангенс \\
\bottomrule
\end{tabular}
\end{center}

\subsection{Элементарный перцептрон}

Элементарный перцептрон вычисляет
\[
  y(x) = \theta\bigl(\langle w, x\rangle\bigr),
  \qquad
  \theta(a) = \begin{cases} 1, & a > 0,\\ -1, & a \le 0. \end{cases}
\]
Геометрически $\langle w, x\rangle = 0$ задаёт гиперплоскость с нормалью $w$, а
перцептрон относит точку к одному из классов по тому, с какой стороны от неё
точка лежит. Свободный член включается в $w$ добавлением к $x$ фиктивной
компоненты, равной единице.

Классы \emph{линейно отделимы}, если существует гиперплоскость, разделяющая их,
то есть существует $\est w$ с
\[
  \langle \est w, x\rangle > 0 \ \ \forall x \in \mathcal{C}_1,
  \qquad
  \langle \est w, x\rangle < 0 \ \ \forall x \in \mathcal{C}_2 .
\]

\subsection{Алгоритм обучения перцептрона}

Обучающая выборка $\mathcal{H} = \mathcal{H}_1 \cup \mathcal{H}_2$ размера $N$,
элементы пронумерованы и перебираются циклически: $x(n) = x(n \bmod N)$.
Начальные веса $w(0) = 0$. На каждой итерации:
\begin{ruslist}
  \item[--] если классификация $x(n)$ верна, веса не меняются: $w(n+1) = w(n)$;
  \item[--] если $x(n) \in \mathcal{H}_1$ классифицирован неверно:
    $w(n+1) = w(n) + \eta\,x(n)$;
  \item[--] если $x(n) \in \mathcal{H}_2$ классифицирован неверно:
    $w(n+1) = w(n) - \eta\,x(n)$.
\end{ruslist}
Изменение весов будем называть \emph{коррекцией}. Алгоритм останавливается, когда
за $N$ итераций подряд, то есть за полный проход по выборке, не было ни одной
коррекции. Алгоритм сходится, если за конечное число итераций веса перестают
меняться и классификация верна для всей выборки.

\subsection{Теорема о сходимости и её доказательство}

\begin{theorem}
Если классы линейно отделимы, алгоритм обучения перцептрона сходится.
\end{theorem}

Схема доказательства такова. Полагается $\eta = 1$, $w(0) = 0$, и
рассматривается случай, когда все встречающиеся $x(k)$ лежат в $\mathcal{H}_1$ и
классифицируются неверно. Тогда после $n$ коррекций
$w(n) = x(0) + x(1) + \dots + x(n-1)$. Вводятся
\[
  \alpha = \min_{x \in \mathcal{H}_1}\langle \est w, x\rangle > 0,
  \qquad
  \beta = \max_{x \in \mathcal{H}_1}\norm{x}^2 .
\]

\emph{Оценка снизу.} Умножим $w(n)$ скалярно на $\est w$. Каждое слагаемое не
меньше $\alpha$, поэтому
\[
  \langle \est w, w(n)\rangle = \sum_{i=0}^{n-1} \langle \est w, x(i)\rangle \ge n\alpha .
\]
Отсюда и из неравенства Коши-Буняковского
$\norm{\est w}^2\norm{w(n)}^2 \ge \langle \est w, w(n)\rangle^2$ следует
\begin{equation}\label{eq:lower}
  \norm{w(n)}^2 \ge \frac{(n\alpha)^2}{\norm{\est w}^2} .
\end{equation}

\emph{Оценка сверху.} Из $w(k+1) = w(k) + x(k)$
\[
  \norm{w(k+1)}^2 = \norm{w(k)}^2 + \norm{x(k)}^2 + 2\langle w(k), x(k)\rangle,
\]
а при неверной классификации $\langle w(k), x(k)\rangle \le 0$, поэтому
$\norm{w(k+1)}^2 - \norm{w(k)}^2 \le \norm{x(k)}^2$. Суммируя по
$k = 0, \dots, n-1$,
\begin{equation}\label{eq:upper}
  \norm{w(n)}^2 \le n\beta .
\end{equation}

\emph{Вывод.} Нижняя оценка растёт как $n^2$, верхняя как $n$, поэтому обе
выполняются только при $n^2\alpha^2/\norm{\est w}^2 \le n\beta$, то есть при
\begin{equation}\label{eq:n0}
  n \le n_0 = \frac{\beta\norm{\est w}^2}{\alpha^2} .
\end{equation}
Больше $n_0$ коррекций быть не может.

Три упрощения этого рассуждения снимаются в заданиях 2 и 3: произвольное $\eta$,
элементы обоих классов, наличие верных классификаций на промежуточных итерациях.

\subsection{Многослойный перцептрон}

Многослойный перцептрон вычисляет последовательность преобразований
\[
  z^{(l)} = \varphi\bigl(W_l z^{(l-1)} + b_l\bigr), \qquad l = 1,\dots,L,
  \qquad z^{(0)} = x,
\]
где $\varphi$ -- функция активации, применяемая покомпонентно, $W_l$ имеет
размер $h_l \times h_{l-1}$, а $h_l$ -- число узлов $l$-го слоя. Активация
\emph{линейна}, если $\varphi(a) = a$ (или, в общем виде,
$\varphi(a) = \kappa a + \lambda$).

Понадобится свойство ранга: $\mathrm{rank}(AB) \le \min(\mathrm{rank}\,A, \mathrm{rank}\,B)$,
а $\mathrm{rank}$ матрицы размера $p \times q$ не превосходит $\min(p,q)$.

\section{Задания}

\zad{1}
\dano Алгоритм обучения перцептрона и доказательство теоремы о сходимости.

\treb Выяснить, важен ли порядок, в котором пронумерованы векторы признаков
обучающей выборки.

\zachem Порядок предъявления выбирается при реализации алгоритма, и нужно знать,
может ли он помешать сходимости.

\resh Разделим два вопроса: влияет ли порядок на \emph{факт} сходимости и на
\emph{результат} работы алгоритма.

\emph{На сходимость не влияет.} В оценку \eqref{eq:n0} входят только $\alpha$,
$\beta$ и $\norm{\est w}$. Величины $\alpha$ и $\beta$ определены как минимум и
максимум по выборке как по \emph{множеству}, а $\est w$ существует по условию
линейной отделимости, тоже независимо от нумерации. Значит и оценка $n_0$, и сам
факт сходимости от порядка не зависят: при любой нумерации коррекций будет не
более $n_0$.

\emph{На результат влияет.} Траектория $w(n)$, фактическое число коррекций и
итоговый вектор весов зависят от того, в каком порядке встретились неверно
классифицированные элементы. Причина в том, что каждая коррекция прибавляет
конкретный $x(n)$, и сумма зависит от того, какие именно элементы успели
оказаться неверно классифицированными. Алгоритм находит \emph{какой-то}
разделяющий вектор из бесконечного множества подходящих, а какой именно --
определяется порядком.

\emph{Одно требование к нумерации всё же есть}: каждый элемент выборки должен
встречаться бесконечно много раз. Циклический перебор $x(n) = x(n \bmod N)$ это
обеспечивает. Если же какой-то элемент перестанет предъявляться, алгоритм может
остановиться на векторе, который на этом элементе ошибается: критерий остановки
проверяет неизменность весов на протяжении $N$ подряд идущих итераций, и смысл
он имеет только при полном переборе выборки за эти $N$ шагов.

\zad{2}
\dano Алгоритм обучения перцептрона с коэффициентом обучения $\eta$.

\treb Показать, что если алгоритм сходится при $\eta = 1$, то он сходится и при
меньших $\eta$, но потребует большего числа итераций.

\zachem Коэффициент обучения задаётся при реализации. Задание показывает, на что
он влияет в алгоритме перцептрона, а на что нет.

\resh Рассмотрим сначала случай $w(0) = 0$, как в доказательстве теоремы.

\emph{Сходимость сохраняется.} Докажем по индукции, что
$w_\eta(n) = \eta\,w_1(n)$, где индекс обозначает коэффициент обучения. База:
$w_\eta(0) = 0 = \eta\cdot 0$. Переход: пусть $w_\eta(n) = \eta w_1(n)$. Проверка
классификации сводится к сравнению $\langle w, x\rangle$ с нулём, а умножение
$w$ на положительное число $\eta$ знака скалярного произведения не меняет:
\[
  \langle w_\eta(n), x(n)\rangle = \eta\,\langle w_1(n), x(n)\rangle .
\]
Значит оба запуска на этой итерации принимают одно и то же решение. Если
коррекции нет, веса не меняются в обоих. Если есть, то
\[
  w_\eta(n+1) = \eta w_1(n) \pm \eta x(n) = \eta\bigl(w_1(n) \pm x(n)\bigr)
              = \eta\,w_1(n+1).
\]
Индукция завершена. Следовательно последовательность решений одна и та же, и
алгоритм сходится при любом $\eta > 0$ ровно за то же число коррекций, отличаясь
только масштабом весов.

Тот же вывод даёт и оценка. После $n$ коррекций $w_\eta(n) = \eta w_1(n)$, и в
\eqref{eq:lower} и \eqref{eq:upper} обе части умножаются на $\eta^2$:
\[
  \frac{(n\eta\alpha)^2}{\norm{\est w}^2} \le \norm{w_\eta(n)}^2 \le n\eta^2\beta
  \;\Longrightarrow\;
  n \le n_0 = \frac{\beta\norm{\est w}^2}{\alpha^2},
\]
и $\eta$ сокращается: оценка от него не зависит.

\emph{Откуда берётся рост числа итераций.} Утверждение про большее число
итераций относится к случаю $w(0) \ne 0$, когда приведённая индукция уже не
проходит: начальный вектор на $\eta$ не масштабируется. Разделим в этом случае
все веса на $\eta$ и обозначим $v(n) = w(n)/\eta$. Правило коррекции превращается
в $v(n+1) = v(n) \pm x(n)$, а проверка знака при делении на положительное $\eta$
не меняется. Значит
\begin{quote}
работа с коэффициентом $\eta$ из начальной точки $w(0)$ совпадает с работой с
коэффициентом $1$ из начальной точки $w(0)/\eta$.
\end{quote}
Чем меньше $\eta$, тем больше норма стартового вектора $w(0)/\eta$. Если этот
вектор классифицирует часть выборки неверно, каждая коррекция сдвигает его лишь
на $\norm{x} \le \sqrt\beta$, поэтому число коррекций, нужных чтобы его
перебить, растёт пропорционально $\norm{w(0)}/\eta$, то есть как $1/\eta$.

Формально это видно из тех же оценок с $w(0) \ne 0$: после $n$ коррекций
$\langle \est w, w(n)\rangle \ge \langle \est w, w(0)\rangle + n\eta\alpha$ и
$\norm{w(n)}^2 \le \norm{w(0)}^2 + n\eta^2\beta$, и неравенство Коши-Буняковского
даёт
\[
  \bigl(\langle \est w, w(0)\rangle + n\eta\alpha\bigr)^2
  \le \norm{\est w}^2\bigl(\norm{w(0)}^2 + n\eta^2\beta\bigr) .
\]
Разделив на $\eta^2$, получаем ту же оценку с заменой $w(0)$ на $w(0)/\eta$:
граница для числа коррекций растёт при уменьшении $\eta$.

\zad{3}
\dano Доказательство теоремы о сходимости проведено в предположении, что
$x(n) \in \mathcal{H}_1$ и классификация производится неверно на всех итерациях
вплоть до $(n+1)$-й.

\treb Показать, что при других возможных промежуточных результатах работы
алгоритма (есть итерации с верной классификацией, встречаются элементы класса
$\mathcal{H}_2$) проведённых рассуждений достаточно для вывода о сходимости.

\zachem Доказательство теоремы проведено в упрощённой ситуации, задание
дополняет его до полного.

\resh Снимем оба ограничения по очереди.

\emph{Элементы обоих классов.} Введём для каждого элемента выборки вектор
\[
  z(i) = \begin{cases}
    x(i), & x(i) \in \mathcal{H}_1,\\
    -x(i), & x(i) \in \mathcal{H}_2 .
  \end{cases}
\]
Проверим, что в новых обозначениях оба случая выглядят одинаково.

Неверная классификация $x(i) \in \mathcal{H}_1$ означает
$\langle w, x(i)\rangle \le 0$, а неверная классификация
$x(i) \in \mathcal{H}_2$ означает $\langle w, x(i)\rangle > 0$. В обоих случаях
$\langle w, z(i)\rangle \le 0$, и только это условие используется в оценке
сверху. Правило коррекции: для $\mathcal{H}_1$ это $w + \eta x(i)$, для
$\mathcal{H}_2$ это $w - \eta x(i)$, то есть в обоих случаях
\[
  w \leftarrow w + \eta\,z(i) .
\]
Линейная отделимость даёт $\langle \est w, z(i)\rangle > 0$ сразу для всех $i$,
поэтому величины
\[
  \alpha = \min_i \langle \est w, z(i)\rangle > 0,
  \qquad
  \beta = \max_i \norm{z(i)}^2 = \max_i \norm{x(i)}^2
\]
определены и положительны. Выкладки, приводящие к \eqref{eq:lower} и
\eqref{eq:upper}, дословно повторяются с заменой $x$ на $z$, поскольку в них
использовались только вид коррекции и знак $\langle \est w, x\rangle$.

Обратим внимание на одну тонкость: отделимость должна быть строгой, то есть
$\langle \est w, x\rangle < 0$ для $x \in \mathcal{C}_2$, а не $\le 0$. При
нестрогом неравенстве возможно $\alpha = 0$, и тогда оценка \eqref{eq:n0} теряет
смысл. Для конечной выборки строгая отделимость следует из нестрогой: если
$\delta = \min_{x \in \mathcal{H}_1}\langle \est w, x\rangle > 0$, то уменьшение
свободного члена $\est w$ на $\delta/2$ сдвигает гиперплоскость и даёт
$\langle \est w, x\rangle \ge \delta/2$ на $\mathcal{H}_1$ и
$\langle \est w, x\rangle \le -\delta/2$ на $\mathcal{H}_2$.

\emph{Верные классификации на промежуточных итерациях.} Такие итерации оставляют
веса без изменения. Значит они не добавляют слагаемых ни в сумму
$\langle \est w, w\rangle = \sum \langle \est w, z(k)\rangle$, ни в оценку
$\norm{w}^2 \le \sum \norm{z(k)}^2$, где суммы идут по сделанным коррекциям.
Поэтому обе оценки остаются верными, если читать $n$ как число
\emph{коррекций}, а не как номер итерации. Именно так они и записаны:
неравенства \eqref{eq:lower} и \eqref{eq:upper} связывают норму весов с числом
уже сделанных коррекций.

\emph{Завершение.} Из \eqref{eq:lower} и \eqref{eq:upper} число коррекций не
превосходит $n_0$. После последней коррекции веса не меняются, а циклический
перебор гарантирует, что за следующие $N$ итераций будут предъявлены все
элементы выборки, и ни на одном из них коррекции не произойдёт. Значит
классификация верна для всей выборки, и критерий остановки срабатывает. Это и
есть определение сходимости. $\blacksquare$

\zad{4}
\dano Обучающая выборка состоит из одного элемента $\{x, +1\}$, то есть
$x \in \mathcal{H}_1$. Начальные веса $w(0)$ выбираются случайно: компоненты
независимы и равномерно распределены на $[0,1]$.

\treb Найти математическое ожидание числа неверных классификаций подряд.
Выяснить, как оно зависит от $\norm{x}$, и показать, что чем меньше коэффициент
обучения, тем оно больше.

\zachem На простейшем примере видно, как число ошибок при обучении зависит от
начальных весов, длины вектора признаков и коэффициента обучения.

\resh Так как элемент один, все итерации предъявляют один и тот же $x$, и
коррекции идут подряд, пока классификация неверна. После $k$ коррекций
\[
  w_k = w(0) + k\eta x,
  \qquad
  \langle w_k, x\rangle = \langle w(0), x\rangle + k\eta\norm{x}^2 .
\]
Обозначим $S = \langle w(0), x\rangle$. Скалярное произведение растёт с каждой
коррекцией на постоянную величину $\eta\norm{x}^2 > 0$, поэтому число неверных
классификаций $K$ -- это наименьшее $k$, при котором оно становится
положительным:
\[
  K = \begin{cases}
    0, & S > 0,\\[1mm]
    \Bigl\lfloor \dfrac{-S}{\eta\norm{x}^2}\Bigr\rfloor + 1, & S \le 0 .
  \end{cases}
\]

Выделим направление: $x = \norm{x}\,u$, где $\norm{u} = 1$, и положим
$t = -\langle w(0), u\rangle$. Тогда $S = -\norm{x}\,t$ и
\[
  \frac{-S}{\eta\norm{x}^2} = \frac{\norm{x}\,t}{\eta\norm{x}^2}
  = \frac{t}{\eta\norm{x}},
\]
поэтому
\begin{equation}\label{eq:K}
  K = \Ind{t \ge 0}\Bigl(\Bigl\lfloor \frac{t}{c}\Bigr\rfloor + 1\Bigr),
  \qquad c = \eta\norm{x} .
\end{equation}

Отсюда сразу видно главное: $K$ зависит от $\eta$ и $\norm{x}$ только через их
произведение $c = \eta\norm{x}$. Величина $t$ от $\eta$ и $\norm{x}$ не зависит,
она определяется только направлением $u$ и распределением $w(0)$.

Оценим математическое ожидание. Из $t/c - 1 < \lfloor t/c\rfloor \le t/c$
следует
\begin{equation}\label{eq:EK}
  \frac{\Exp t_+}{c} \;\le\; \Exp K \;\le\; \frac{\Exp t_+}{c} + \PP(t \ge 0),
  \qquad t_+ = \max(t, 0).
\end{equation}
Отсюда асимптотика: $\Exp K \to \PP(t \ge 0)$ при $c \to \infty$ и
$\Exp K \sim \Exp t_+/c$ при $c \to 0$.

\emph{Зависимость от $\norm{x}$.} По \eqref{eq:EK} при малых $c$ величина
$\Exp K$ обратно пропорциональна $\norm{x}$: чем длиннее вектор, тем сильнее
каждая коррекция сдвигает скалярное произведение, и тем быстрее оно становится
положительным.

\emph{Зависимость от $\eta$.} При каждом значении $t$ величина $K$ из
\eqref{eq:K} не возрастает с ростом $c = \eta\norm{x}$, поэтому $\Exp K$ не
убывает при уменьшении $\eta$. При малых $c$ по \eqref{eq:EK} она обратно
пропорциональна $\eta$: уменьшение коэффициента обучения вдвое примерно вдвое
увеличивает число ошибок подряд. Это и требовалось показать. $\blacksquare$

\emph{Частный случай.} Если $x \ne 0$ и все его компоненты неотрицательны
(например $x$ -- яркости пикселей), то $S = \langle w(0), x\rangle > 0$ почти наверное, так
как все компоненты $w(0)$ положительны почти наверное. Тогда $K = 0$: ни одной
ошибки не произойдёт независимо от $\eta$.

\zad{5}
\dano Нейронная сеть, в которой функция активации линейна.

\treb Показать, что задаваемое такой сетью преобразование линейно.

\zachem Объясняет, зачем в нейронных сетях нужны нелинейные функции активации.

\resh При $\varphi(a) = a$ слой вычисляет $z^{(l)} = W_l z^{(l-1)} + b_l$.
Развернём композицию слоёв, начиная с $z^{(0)} = x$:
\[
  z^{(1)} = W_1 x + b_1,
  \qquad
  z^{(2)} = W_2 W_1 x + W_2 b_1 + b_2,
\]
и по индукции
\[
  z^{(L)} = \underbrace{W_L W_{L-1}\cdots W_1}_{M}\,x
          + \underbrace{\sum_{l=1}^{L} W_L \cdots W_{l+1}\,b_l}_{c}
          = Mx + c .
\]
Это аффинное преобразование, при $b_l = 0$ линейное. Если активация имеет общий
линейный вид $\varphi(a) = \kappa a + \lambda$, то множитель $\kappa$ вносится в
$W_l$, а сдвиг $\lambda$ в $b_l$, и результат тот же. $\blacksquare$

Следствие: любое число слоёв с линейной активацией сворачивается в один слой с
матрицей $M$, поэтому добавление слоёв имеет смысл только при нелинейной
активации.

\zad{6}
\dano Сеть с линейными функциями активации, в которой количество узлов
некоторого скрытого слоя меньше числа входов или числа выходов сети.

\treb Показать, что задаваемое сетью линейное преобразование имеет меньшую
размерность, чем размерность входа или выхода.

\zachem Узкий скрытый слой используется в автоэнкодерах. Задание показывает, что
он ограничивает передачу информации со входа на выход при любых весах.

\resh По заданию 5 сеть задаёт преобразование $x \mapsto Mx + c$ с матрицей
$M = W_L\cdots W_1$, где $W_l$ имеет размер $h_l \times h_{l-1}$. Размерность
преобразования есть размерность его образа, то есть $\mathrm{rank}\,M$. Ранг
произведения не превосходит ранга каждого сомножителя, а ранг матрицы не
превосходит меньшей из её размерностей, поэтому
\[
  \mathrm{rank}\,M \le \min_{1 \le l \le L}\mathrm{rank}\,W_l
                  \le \min_{0 \le l \le L} h_l .
\]
Пусть скрытый слой содержит $h$ узлов, тогда $\mathrm{rank}\,M \le h$.

Если $h < h_0$ (меньше входов), то $\mathrm{rank}\,M < h_0$, и ядро нетривиально:
$\dim\ker M = h_0 - \mathrm{rank}\,M \ge h_0 - h > 0$. Входы, разность которых
лежит в ядре, дают одинаковый выход.

Если $h < h_L$ (меньше выходов), то $\mathrm{rank}\,M < h_L$: образ лежит в
подпространстве размерности не выше $h$, и не всякий выход достижим.
$\blacksquare$

\zad{7}
\dano Нейронная сеть с линейными функциями активации, в которой допустимы дуги,
соединяющие узлы из разных слоёв, а не только соседних.

\treb Показать, что для ограниченных входных значений такая сеть эквивалентна
многослойному перцептрону с линейными скрытыми узлами.

\zachem Показывает, что соединения через слой, как в остаточных сетях, при
линейных активациях не расширяют класс задаваемых сетью преобразований.

\resh Докажем два включения.

\emph{Такая сеть задаёт аффинное преобразование.} Граф сети ацикличен, поэтому
его узлы можно упорядочить топологически: каждая дуга ведёт от более раннего
узла к более позднему. Идём по этому порядку. Значение входного узла есть $x_i$,
то есть аффинная функция от $x$. Значение любого другого узла есть линейная
комбинация значений его предшественников плюс сдвиг, а линейная комбинация
аффинных функций аффинна. По индукции значение каждого узла, включая выходные,
есть аффинная функция от $x$.

\emph{Любое такое преобразование реализуется обычным МСП.} Достаточно показать,
как убрать дуги, идущие через слой. Расширим состояние каждого слоя копией всех
предыдущих:
\[
  \hat z^{(l)} = \bigl(z^{(l)},\, z^{(l-1)},\, \dots,\, z^{(0)}\bigr).
\]
Тогда $\hat z^{(l)} = M_l\,\hat z^{(l-1)}$, где первая блочная строка $M_l$
содержит веса всех дуг, входящих в слой $l$ (в том числе идущих из далёких
слоёв, значения которых теперь доступны в $\hat z^{(l-1)}$), а остальные блоки
единичные и просто пробрасывают значения дальше. Все узлы линейные, все дуги
соединяют только соседние слои.

Пример для дуги, идущей со входа сразу на выход,
$y = W_2 W_1 x + S x$:
\[
  M_1 = \begin{pmatrix} W_1 \\ I \end{pmatrix},
  \qquad
  M_2 = \begin{pmatrix} W_2 & S \end{pmatrix},
  \qquad
  M_2 M_1 x = W_2 W_1 x + S x = y . \qquad\blacksquare
\]

Замечание про ограниченность входа. Для линейной активации равенство точное, и
никакой ограниченности не требуется. Условие становится существенным, если
линейность активации сама является приближением на ограниченном участке: так,
$\sigma(a) \approx 1/2 + a/4$ вблизи нуля, и рассуждение работает лишь пока
предактивации остаются малыми.

\zad{8}
\dano Логистический сигмоид и гиперболический тангенс
\[
  \sigma(a) = \frac{1}{1 + e^{-a}},
  \qquad
  \tanh(a) = \frac{e^a - e^{-a}}{e^a + e^{-a}} .
\]

\treb Найти их производные и объяснить, почему такие функции активации удобны
при реализации метода обратного распространения ошибки.

\zachem Производные функций активации вычисляются на каждом шаге обратного
распространения ошибки.

\resh Дифференцируем сигмоид как сложную функцию и группируем множители:
\[
  \sigma'(a) = \frac{e^{-a}}{(1+e^{-a})^2}
             = \frac{1}{1+e^{-a}}\cdot\frac{e^{-a}}{1+e^{-a}}
             = \sigma(a)\bigl(1 - \sigma(a)\bigr),
\]
где во второй дроби использовано
$\dfrac{e^{-a}}{1+e^{-a}} = 1 - \dfrac{1}{1+e^{-a}}$.

Для гиперболического тангенса дифференцируем как частное:
\[
  \tanh'(a) = \frac{(e^a+e^{-a})^2 - (e^a-e^{-a})^2}{(e^a+e^{-a})^2}
            = 1 - \tanh^2(a).
\]

\emph{Почему это удобно.} Метод обратного распространения устроен так: сначала
прямой проход вычисляет активации всех узлов, затем обратный проход перемножает
по цепному правилу производные активаций. Обе найденные производные выражаются
\emph{через само значение функции}, а не через её аргумент. Поэтому прямой
проход достаточно сохранить в виде значений $z = \sigma(a)$, и на обратном
проходе производная получается как $z(1-z)$: не нужно ни хранить предактивации
$a$, ни заново вычислять экспоненты. Вместо вызова $\exp$ на каждый узел
остаётся одно умножение, что при миллионах узлов существенно.

Дополнительно обе функции монотонны, гладки на всей прямой и имеют ограниченные
производные ($\sigma' \le 1/4$, $\tanh' \le 1$), поэтому множители в цепном
правиле не растут. Обратная сторона той же ограниченности -- затухание
градиента: при большом числе слоёв произведение множителей, меньших единицы,
стремится к нулю.

Функции связаны соотношением $\tanh(a) = 2\sigma(2a) - 1$, то есть с точностью
до масштаба это одна и та же функция; $\tanh$ отличается тем, что центрирован
относительно нуля, а его производная в нуле равна 1 против $1/4$ у сигмоида.

\end{document}
